\chapter{程序开发}

本章从程序结构角度说明领域实体、分类策略、账单解析和报表生成的设计，并进一步介绍数据库、模块间数据契约、设计模式及核心业务流程。重点在于说明设计职责和协作关系，而非逐项解释源代码。

\section{面向对象程序组织}

系统根据问题特征选择程序组织方式：领域对象、可替换分类策略、平台解析器和报表构建过程采用面向对象设计；身份校验、数据聚合等流程则以分层服务组织。二者共同服务于职责分离、扩展性和可测试性。

\subsection{领域实体及其关联关系}

系统使用对象关系映射（ORM）描述 12 个领域实体，并采用数据映射模式将领域结构与数据库访问分离。实体关系围绕五个业务域展开：用户拥有个人交易与自定义分类；导入批次包含一个或多个上传文件，并产生标准化交易；交易保留自动分类建议、人工最终分类和历次分类记录；家庭组织通过成员关系连接多个用户；报表任务记录生成过程并关联最终文件。

交易去重依据发生时间、金额、商户、商品和备注等规范化信息生成稳定摘要，并由数据库唯一约束阻止同一用户重复保存相同交易。分类设计同时保留“自动建议”和“人工最终决定”，各项分析优先采用人工确认结果。分类历史采用追加方式保存，使分类来源和结果变化可以追溯。图~\ref{fig:entity-class} 仅列出参与领域关系的关键属性，其他实现字段由数据库设计图补充。

\begin{figure}[p]
  \centering
  \includegraphics[
    width=0.92\textwidth,
    height=0.76\textheight,
    keepaspectratio
  ]{entity-domain-model.png}
  \caption{领域实体类图}
  \label{fig:entity-class}
\end{figure}

\subsection{分类器策略}

分类子系统采用策略、责任链和外观三种模式。缓存分类、模型分类和规则分类遵循统一的数据契约，可以根据运行配置替换；分类请求依次尝试已有缓存、模型服务和规则后备方案，在外部服务不可用时仍能给出可解释结果；统一入口负责协调缓存、模型调用、失败重试和结果保存，使导入模块不需要了解各类分类方式的内部差异。

外部模型调用由等待队列控制，并采用指数退避减少持续失败时的请求压力。无论结果来自缓存、外部模型、本地模型还是规则，系统都统一记录类别、置信度、理由和来源，便于后续展示、人工校正和质量分析。分类器之间的静态关系如图~\ref{fig:classifier-class} 所示。

\begin{figure}[!htbp]
  \centering
  \includegraphics[
    width=0.66\textwidth,
    height=0.30\textheight,
    keepaspectratio
  ]{分类器策略类图.png}
  \caption{分类器策略类图}
  \label{fig:classifier-class}
\end{figure}

\subsection{账单解析器}

解析子系统负责将微信和支付宝的异构账单转换为统一交易结构，采用策略、适配器与注册表思想。系统首先根据文件名和内容特征识别平台，再处理不同编码、文件格式和表头位置，最后将平台特有列名映射为统一的时间、金额、商户、商品和状态等信息。规范化阶段统一字符形式并生成去重摘要。

不同平台的解析规则彼此独立，新平台只需提供识别规则和字段映射，不需要修改分类、分析和报表模块。统一读取能力屏蔽 CSV 与 Excel、字符编码和表头位置等差异；注册机制负责选择合适的平台解析策略。该结构降低了新增账单来源对其他模块的影响，类之间的关系如图~\ref{fig:parser-class} 所示。

\begin{figure}[H]
  \centering
  \includegraphics[width=0.92\textwidth]{解析器类图.png}
  \caption{账单解析器类图}
  \label{fig:parser-class}
\end{figure}

\subsection{报表构建器}

报表子系统采用构建器模式，将数据准备、封面生成、摘要指标、分类汇总、商户排行和交易明细等步骤有序组合。数据聚合与页面渲染相互分离，同一份报表数据可以被多个页面区域复用；各渲染步骤职责单一，增加新章节时只需扩展构建过程，而不必修改已有统计逻辑。

系统同时考虑中文字体兼容和输出文件记录，使报表在不同运行环境下保持可读，并能够通过任务记录追踪生成结果。报表构建器的结构如图~\ref{fig:reportbuilder-class} 所示。

\begin{figure}[H]
  \centering
  \includegraphics[width=0.85\textwidth]{报表构建器类图.png}
  \caption{报表构建器类图}
  \label{fig:reportbuilder-class}
\end{figure}

\section{数据库设计}

系统数据库共 12 张表，按业务域分两组展示。

\subsection{核心账单与分类}

核心域由用户、分类、分类规则、导入批次、上传文件、交易、分类历史和分类缓存八类数据组成。

\begin{figure}[H]
  \centering
  \includegraphics[
    width=0.60\textwidth,
    height=0.56\textheight,
    keepaspectratio
  ]{核心账单ER图.png}
  \caption{核心账单与分类 ER 图}
  \label{fig:er-core}
\end{figure}

关键设计决策包括：以用户和交易摘要的联合唯一约束阻止重复导入；分别保存自动分类建议与人工最终分类，并优先采用人工结果；以复核标记识别低置信度或未分类交易；分类历史只追加不覆盖，以保留每次分类记录；分类缓存按用户、分类服务和文本摘要隔离，避免不同用户之间共享敏感结果。

\subsection{家庭组织与报表扩展}

扩展域由家庭组织、组织成员、报表任务和生成报表四类数据组成，并与用户和交易数据共享关联。

组织成员角色分为所有者、管理员和普通成员：所有者拥有完整管理权限，管理员可管理普通成员，普通成员仅可查看授权的聚合数据。家庭组织不改变交易的个人归属，家庭分析在查询时依据成员关系扩展数据范围。报表任务记录生成请求的处理阶段，生成报表记录最终文件及其所属用户。

\begin{figure}[H]
  \centering
  \includegraphics[width=0.92\textwidth]{家庭组织与报表扩展ER图.png}
  \caption{家庭组织与报表扩展 ER 图}
  \label{fig:er-family}
\end{figure}

\section{接口设计}

\subsection{外部 REST API}

系统对外接口采用统一前缀管理，除注册、登录与健康检查外均要求身份凭证。表~\ref{tab:api-overview} 集中列出主要接口及访问权限；更详细的请求与响应结构由 OpenAPI 接口文档描述。

\begingroup
\footnotesize
\setlength{\tabcolsep}{3pt}
\begin{longtable}{@{}p{2.3cm} p{3.2cm} p{2.5cm} p{2.1cm} p{2.2cm}@{}}
\caption{主要外部 API 接口}
\label{tab:api-overview}\\
\toprule
\textbf{方法} & \textbf{路径} & \textbf{功能} & \textbf{权限} & \textbf{关键输入/输出} \\
\midrule
\endfirsthead
\multicolumn{5}{c}{\small 表~\thetable\quad 主要外部 API 接口（续）} \\
\toprule
\textbf{方法} & \textbf{路径} & \textbf{功能} & \textbf{权限} & \textbf{关键输入/输出} \\
\midrule
\endhead
\midrule
\multicolumn{5}{r}{\small 接下页} \\
\endfoot
\bottomrule
\endlastfoot
POST & \path{/api/auth/register}、\path{/login} & 注册或登录并返回身份凭证 & 无 & 用户名、密码、用户信息与身份凭证 \\
\hline
GET & \path{/api/auth/me} & 获取当前用户 & 登录用户 & 当前用户信息 \\
\hline
\shortstack[l]{GET/POST/\\PATCH} & \path{/api/categories} & 查询、创建和更新分类 & 登录用户 & 分类信息 \\
\hline
POST & \path{/api/imports} & 上传一个或多个账单文件 & 登录用户 & 文件与导入批次信息 \\
\hline
\shortstack[l]{GET/\\DELETE} & \path{/api/imports[/{batch_id}]} & 查询批次、文件或删除本人批次 & 资源所有者 & 批次状态、进度、错误信息 \\
\hline
GET & \path{/api/transactions} & 分页筛选交易 & 登录用户 & 日期、平台、分类、文件、关键词 \\
\hline
PATCH & \path{/api/transactions/{id}/category} & 确认或修改最终分类 & 资源所有者 & \texttt{category\_id}、\texttt{mark\_reviewed} \\
\hline
POST & \path{/api/classification/reclassify} & 对本人指定交易重新分类 & 登录用户 & 交易编号列表、分类服务 \\
\hline
\shortstack[l]{GET/\\POST} & \path{/api/classification/retry-status}、\path{/retry-all} & 查询重试状态或恢复失败交易 & 系统管理员 & 聚合状态或恢复数量 \\
\hline
GET & \path{/api/dashboard/summary} & 获取个人或家庭聚合分析 & 登录用户/组织成员 & 日期、文件、组织编号 \\
\hline
\shortstack[l]{GET/\\POST} & \path{/api/personality/profile}、\path{/quiz/result} & 获取画像或提交问卷 & 登录用户/组织成员 & 画像、健康评分、偏差分析 \\
\hline
\shortstack[l]{GET/POST\\PATCH/DELETE} & \path{/api/organizations} & 组织与成员管理 & 登录用户；写操作按角色 & 组织、成员和角色信息 \\
\hline
\shortstack[l]{POST/\\GET} & \path{/api/reports}、\path{/reports/{id}/download} & 生成或下载 PDF 报表 & 登录用户/组织成员、所有者 & 日期、文件、报表文件 \\
\hline
\shortstack[l]{GET/\\PUT} & \path{/api/settings} & 查询或保存用户分类设置 & 登录用户 & 分类服务、模型名、置信度阈值 \\
\hline
GET & \path{/health} & 应用存活探测 & 无 & 服务状态 \\
\end{longtable}
\endgroup

接口错误遵循 HTTP 状态码语义：未提供有效凭据返回 401，组织角色不足返回 403，资源不存在或不属于当前用户时通常返回 404，请求字段不符合 Pydantic 模型时返回 422，解析或文件生成等内部异常返回 500。资源归属场景使用 404 可以避免向调用者泄露目标资源是否存在。

\subsection{模块间数据契约}

\begin{table}[H]
\centering
\small
\caption{模块间数据契约}
\label{tab:internal-contracts}
\begin{tabularx}{0.98\textwidth}{p{3.3cm} p{4.2cm} X}
\toprule
\textbf{协作边界} & \textbf{输入与输出} & \textbf{作用} \\
\midrule
账单解析 & 原始账单文件 $\rightarrow$ 平台信息与标准交易列表 & 隔离微信、支付宝等原始列名差异，后续模块只使用统一交易结构。 \\
\hline
智能分类 & 标准交易 $\rightarrow$ 分类结果 & 协调缓存、模型、规则后备方案和失败重试，调用方不依赖具体分类方式。 \\
\hline
分类结果 & 类别、子类别、置信度、理由、来源 & 统一规则、外部服务和本地模型的输出格式，便于保存与展示。 \\
\hline
家庭成员范围 & 组织编号 $\rightarrow$ 可访问成员范围 & 将家庭权限校验与财务分析、消费人格、报表聚合相互分离。 \\
\hline
报表生成 & 用户范围、日期和文件筛选 $\rightarrow$ PDF 报表 & 将报表数据准备、页面渲染和结果记录组织为统一过程。 \\
\hline
接口数据模型 & JSON 或表单输入 $\leftrightarrow$ 数据传输对象 & 在 HTTP 边界执行字段校验和序列化，避免直接暴露数据库实体。 \\
\bottomrule
\end{tabularx}
\end{table}

\subsection{任务与状态语义}

\begin{table}[H]
\centering
\small
\caption{任务状态与故障语义}
\label{tab:task-status}
\begin{tabularx}{0.98\textwidth}{p{3cm} p{4.2cm} X}
\toprule
\textbf{任务} & \textbf{状态或流转} & \textbf{说明} \\
\midrule
账单导入 & 等待处理 $\rightarrow$ 处理中 $\rightarrow$ 完成/部分失败/失败 & 后台处理过程更新总数与已处理数；部分记录失败时保留成功数据并记录原因。 \\
\hline
外部模型分类 & 等待重试 $\rightarrow$ 分类完成/重试失败 & 后台任务处理器依次处理等待交易，并记录最终分类来源或失败结果。 \\
\hline
报表生成 & 当前实现：处理中 $\rightarrow$ 完成 & 当前采用同步方式构建 PDF；系统已经预留等待和失败状态，但异步提交、状态查询和失败反馈尚未形成完整流程。 \\
\bottomrule
\end{tabularx}
\end{table}

\section{设计模式应用}

下表不只列出模式名称，还说明它解决的变化点以及支撑的质量要求。

\begingroup
\footnotesize
\setlength{\tabcolsep}{3pt}
\begin{longtable}{@{}p{2.4cm} p{4.2cm} p{3.4cm} p{2.0cm}@{}}
    \caption{设计模式应用汇总}
    \label{tab:design-patterns}\\
    \toprule
    \textbf{设计模式} & \textbf{解决的问题} & \textbf{主要应用} & \textbf{支撑属性} \\
    \midrule
    \endfirsthead
    \multicolumn{4}{c}{\footnotesize 表~\thetable\quad 设计模式应用汇总（续）} \\
    \toprule
    \textbf{设计模式} & \textbf{解决的问题} & \textbf{主要应用} & \textbf{支撑属性} \\
    \midrule
    \endhead
    \midrule
    \multicolumn{4}{r}{\footnotesize 接下页} \\
    \endfoot
    \bottomrule
    \endlastfoot
    Data Mapper & 将领域实体与具体数据库查询分离，减少持久化方式对业务结构的影响 & 领域实体与数据访问 & 可维护性 \\
    \hline
    Strategy & 在统一协作方式下替换规则、外部模型或本地模型分类，也可替换不同平台解析规则 & 分类与账单解析 & 可扩展性 \\
    \hline
    Chain of Responsibility & 依次尝试缓存、模型和规则，在前一环节不可用时继续处理 & 复合分类过程 & 可靠性 \\
    \hline
    Facade & 对导入模块隐藏缓存、模型调用、失败重试与结果记录的复杂性 & 分类统一入口 & 可维护性 \\
    \hline
    Adapter/Factory & 将不同文件格式和平台字段转换为统一交易结构，并自动选择合适的解析策略 & 账单读取与平台识别 & 兼容性、可扩展性 \\
    \hline
    Builder & 将数据准备、摘要、图表和明细等报表步骤按序组合 & PDF 报表生成 & 可维护性 \\
    \hline
    Producer-Consumer & 将产生分类任务的业务过程与执行外部模型调用的后台过程分离 & 外部模型失败重试 & 可靠性 \\
\end{longtable}
\endgroup

\section{核心流程设计}

\subsection{账单导入与大模型初分}

用户上传账单后的处理过程包含六个阶段。(1) 系统接收文件、创建导入批次并向用户反馈进度；(2) 识别账单平台，将原始列转换为统一交易结构；(3) 规范化文本并依据交易特征执行去重；(4) 按“缓存、模型服务、规则后备方案”的顺序获取分类结果，外部服务失败时由后台重试机制继续处理；(5) 同时保存当前分类结果与可追溯的分类历史，低置信度交易进入人工复核；(6) 汇总成功与失败记录，形成完成、部分失败或失败的批次结果。

\begin{figure}[H]
  \centering
  \includegraphics[width=0.92\textwidth]{账单导入与大模型初分顺序图.png}
  \caption{账单导入与大模型初分顺序图}
  \label{fig:seq-import}
\end{figure}

\subsection{家庭组织聚合查询}

家庭组织视图的数据聚合不改变交易的个人归属，而是在查询阶段根据成员身份扩展可访问范围。系统首先取得用户所在组织及其角色；当用户选择家庭财务看板、消费人格或报表时，系统校验成员身份并汇总组织内成员的交易数据。财务看板返回家庭收入、支出、净流入与分类分布，消费人格模块计算家庭整体画像与财务健康评分，报表模块据此生成家庭聚合 PDF。

\begin{figure}[!htbp]
  \centering
  \includegraphics[
    width=0.58\textwidth,
    height=0.32\textheight,
    keepaspectratio
  ]{对象图.png}
  \caption{家庭聚合场景对象图}
  \label{fig:obj-family}
\end{figure}
